\section{Introduction}

% Starte generelt med å utforske hva den faglige diskusjonen går i for tiden
% Gå over til mer spesifikt, hva har andre gjort med problemstillingen som utforskes
% Presenter det vitenskapelige målet og problemstillingen, hva er det de andre ikke har fått til som vi skal gjøre

The study of phase transitions and collective behavior in many-particle systems is a central topic in statistical physics. Such systems can exhibit dramatic changes in macroscopic properties as external parameters, such as temperature, are varied. A classic example is ferromagnetism, where a material transitions from a disordered state at high temperature to an ordered state with nonzero magnetization below a critical temperature.
\\ \\
The Ising model provides one of the simplest theoretical frameworks for studying such phenomena. In its two-dimensional form, the model consists of spins 
$ s_i = \pm 1$
arranged on a square lattice, where each spin interacts with its nearest neighbors. Despite its simplicity, the model captures essential features of real magnetic systems, including spontaneous magnetization and critical behavior near the phase transition. The exact solution for the two-dimensional Ising model without an external magnetic field predicts a critical temperature 
$T_c \approx 2.27 J/k_B$
, marking a second-order phase transition between ordered and disordered phases.
\\ \\
In this project, the two-dimensional Ising model is investigated using Monte Carlo simulations based on the Metropolis algorithm. The aim is to study the thermodynamic behavior of the system and characterize the phase transition by analyzing quantities such as energy, magnetization, specific heat, and magnetic susceptibility. The dependence of these observables on temperature is examined in order to estimate the critical temperature and explore critical fluctuations near the transition point.
\\ \\
Additionally, the effects of system size, impurities, and external magnetic fields are considered. Impurities are introduced as fixed spins to mimic disorder in real materials, while simulated annealing is used to explore low-energy configurations. The influence of an external magnetic field is studied through the behavior of magnetization and susceptibility, and scaling relations are tested near the critical region.
\\ \\
Through these simulations, the project demonstrates how numerical methods can reproduce key features of phase transitions and provides insight into the physical behavior of interacting spin systems.
